\section{Literature and Economic Mechanism}

\subsection{Momentum, crashes, and volatility}

Cross-sectional momentum buys past winners and sells past losers
\citep{jegadeesh1993}. Its long-run average profitability coexists with severe episodic
losses. \citet{daniel2016} show that momentum crashes are predictable in panic states
and are associated with market rebounds and time-varying exposure. Volatility management
can reduce such losses, although scaling changes leverage and requires a target-risk
choice \citep{barroso2015}.

\citet{butt2023} replace continuous scaling with discrete strategy selection. Their
switch is motivated by two conditional relations: high market volatility predicts weaker
momentum and stronger short-term reversal. The source paper interprets reversal partly as
compensation for supplying liquidity when demand for immediacy is high. This view is
consistent with evidence that short-horizon return reversals become stronger when market
liquidity is scarce \citep{nagel2012,jegadeesh1990,lehmann1990}.

The source mechanism is therefore richer than the statement ``high volatility causes
reversal.'' Volatility proxies for a joint state containing distress, prior losses,
funding pressure, liquidity demand, changing winner--loser betas, and a possible market
rebound. The router works only if that joint state is sufficiently stable across time and
markets.

\subsection{Why futures are a useful external test}

Futures markets contain well-documented momentum-related phenomena, including time-series
momentum \citep{moskowitz2012} and momentum combined with term-structure information in
commodity allocation \citep{fuertes2010}. These results do not imply that an equity
cross-sectional state variable transfers mechanically. Futures prices respond directly
to inventories, storage economics, production constraints, hedging demand, interest
rates, currencies, and policy. A shock can represent slow information incorporation and
trend continuation rather than temporary liquidity pressure.

This difference suggests a transportability hypothesis. In equities, a market-level
volatility scalar may summarize a relatively coherent distress state. In a broad futures
cross-section, the same scalar may aggregate economically distinct shocks. The proper
test is therefore not only whether the primary router has a higher sample return, but
whether the conditional sleeve ranking is monotonic, representative across high-volatility
months, and stable to a reasonable definition of the futures ``market.''
